% Editable answer key for the unchanged questions in e33.pdf.
% The original e32.pdf is preserved.
\documentclass[11pt,letterpaper]{article}
\usepackage[T1]{fontenc}
\usepackage[margin=0.75in]{geometry}
\usepackage{amsmath,amssymb}
\usepackage{enumitem,array,booktabs}
\usepackage[hidelinks]{hyperref}
\setlength{\parindent}{0pt}
\setlength{\parskip}{3pt}
\setlength{\emergencystretch}{1.5em}
\setlist[enumerate]{leftmargin=2em,itemsep=5pt,topsep=3pt,parsep=0pt}
\newcommand{\curl}{\nabla\times}
\newcommand{\dv}{\nabla\cdot}
\renewcommand{\vec}[1]{\mathbf{#1}}
\newcommand{\dd}{\mathop{}\!\mathrm{d}}
\begin{document}

{\large\bfseries Exam 3 Practice Problems --- Answers}
\par\smallskip
 Equivalent mathematical forms
are valid. In planar flux questions, the normal is outward.
\setlength{\abovedisplayskip}{4pt}
\setlength{\belowdisplayskip}{4pt}
\setlength{\abovedisplayshortskip}{3pt}
\setlength{\belowdisplayshortskip}{3pt}

\begin{enumerate}
\item
\begin{enumerate}[label=(\alph*),leftmargin=1.6em,itemsep=5pt]
\item $\curl\vec F=(8y+3z^2,0,3-x^2)$,
$\dv\vec F=2xy$. Not conservative.
\item $\displaystyle\curl\vec F=
\left(\frac{x^3y^2}{z^2},4z-7,\frac{3x^2y^2}{z}\right)$,
$\displaystyle\dv\vec F=2+\frac{2x^3y}{z}$.
Not conservative on either half-space $z>0$ or $z<0$.
\item $\curl\vec F=(0,3x^2,0)$,
$\dv\vec F=14-2y$. Not conservative.
\item $\curl\vec F=(0,0,0)$,
$\dv\vec F=-xz\cos y$. Conservative, with potential
$f=xz\cos y+C$.
\end{enumerate}

\item Conservative: $\curl\vec F=\vec0$ and
$\boxed{f=x^2y^2+xz+C}$.
Indeed, $P_y=Q_x=4xy$, $P_z=R_x=1$, and $Q_z=R_y=0$.

\item Conservative, with $\boxed{f=xyz+e^x+e^y+e^z+C}$.

\item $\vec F=\nabla(xyz^2)$, so the integral is
$f(1,1,1)-f(0,0,0)=\boxed{1}$.

\item $\boxed{3e-2}$, using the potential from Problem 3.

\item $Q_x-P_y=-1-2y$, so
\[
\int_0^1\int_0^x(-1-2y)\,\dd y\,\dd x=\boxed{-\frac56}.
\]

\item $P_x+Q_y=0$, so the outward flux is $\boxed{0}$.

\item $\boxed{-8\pi}$.
\item $\boxed{8\pi}$.
\item $\boxed{1}$.
\item $\boxed{8}$.

\item $Q_x-P_y=y$, so
\[
\int_0^1\int_{x^2}^{x}y\,\dd y\,\dd x
=\frac12\int_0^1(x^2-x^4)\,\dd x
=\boxed{\frac1{15}}.
\]

\item $P_x+Q_y=3x$, so
\[
\int_0^1\int_{x^2}^{x}3x\,\dd y\,\dd x=\boxed{\frac14}.
\]
\end{enumerate}

\newpage
\begin{enumerate}[start=14]
\item
\[
\vec r_u\times\vec r_v=(2u^2+2v^2,v-u,-2(u+v)),
\]
\[
\boxed{A=\int_0^2\int_0^1
\sqrt{(2u^2+2v^2)^2+(v-u)^2+4(u+v)^2}\,\dd u\,\dd v}.
\]

\item
\[
\vec r_u\times\vec r_v=(v^2-2uv,u^2-2uv,3u^2v^2),
\]
\[
\boxed{m=\int_0^1\int_0^1 e^{u+v}
\sqrt{(v^2-2uv)^2+(u^2-2uv)^2+9u^4v^4}\,\dd u\,\dd v}.
\]

\item
\[
\boxed{A=2\pi\int_0^1e^{5x}\sqrt{1+25e^{10x}}\,\dd x}.
\]

\item Interpreting the scalar function on the spatial surface as
$f(x,y,z)=y^5$,
\[
\boxed{\iint_S f\,\dd\sigma
=2\pi\int_1^4x^{-3/2}\sqrt{1+\frac1{4x^3}}\,\dd x}.
\]
The radius is $x$, the density is $(x^{-1/2})^5$, and the generating
curve has speed $\sqrt{1+1/(4x^3)}$.

\item $\boxed{\pi/3}$.

\item Use $\vec r(r,\theta)=(r\cos\theta,r\sin\theta,r)$,
with $2\le r\le4$, $0\le\theta\le2\pi$. Then
\[
\iint_S z^2\,\dd\sigma
=\int_0^{2\pi}\int_2^4r^2\sqrt2\,r\,\dd r\,\dd\theta
=\boxed{120\sqrt2\,\pi}.
\]

\item
\[
A=\int_0^{2\pi}\int_0^1r\sqrt{1+4r^2}\,\dd r\,\dd\theta
=\boxed{\frac{\pi(5\sqrt5-1)}6}.
\]

\item The spherical band has $\pi/4\le\phi\le3\pi/4$, so
\[
A=\int_0^{2\pi}\int_{\pi/4}^{3\pi/4}\sin\phi\,\dd\phi\,\dd\theta
=\boxed{2\sqrt2\,\pi}.
\]

\item $\boxed{3\sqrt{14}}$.
\end{enumerate}

\newpage
\begin{enumerate}[start=23]
\item With $\vec r(\phi,\theta)=(\sin\phi\cos\theta,
\sin\phi\sin\theta,\cos\phi)$, the outward vector area element is
$\sin\phi\,\vec r\,\dd\phi\,\dd\theta$. The complete setup is
\[
\boxed{\int_0^{2\pi}\int_0^\pi
\bigl(\sin^3\phi\cos^2\theta\sin\theta
+2\sin\phi\sin\theta-\cos^2\phi\bigr)
\sin\phi\,\dd\phi\,\dd\theta}.
\]

\item The downward vector area element is
$(r\cos\theta,r\sin\theta,-r)\,\dd r\,\dd\theta$, giving
\[
\int_0^{2\pi}\int_2^4(r^2\sin\theta\cos\theta-r)\,\dd r\,\dd\theta
=\boxed{-12\pi}.
\]

\item The upward graph normal factor is $(1/3,1,1)$, and
$(3,1,-2)\cdot(1/3,1,1)=0$. Thus the flux is $\boxed{0}$.

\item The plane is $z=6-3x/5-2y/5$, and the upward flux integrand is
$(3y+7x)/5$ over the triangle with vertices $(0,0),(10,0),(0,15)$:
\[
\int_0^{10}\int_0^{15-3x/2}\frac{3y+7x}{5}\,\dd y\,\dd x
=\boxed{575}.
\]

\item With $z_x=2xy^2+1$ and $z_y=2x^2y+1$, the setup is
\[
\boxed{\int_0^1\int_0^1(-xy^2-4x^2y-3)\,\dd x\,\dd y}.
\]
Its value, for a check, is $-23/6$.

\item With $z_x=4x^3$ and $z_y=2y$, the setup is
\[
\boxed{\int_{-1}^{1}\int_{-1}^{1}
\bigl(-4x^5-2y(x^4+y^2)+(x^4+y^2)^2\bigr)\,\dd x\,\dd y}.
\]
Its value, for a check, is $16/9$.

\item $\boxed{-\pi}$. The intended surface includes its equatorial rim
($z\ge0$); the rim is the boundary used in Stokes' theorem.
\item $\boxed{0}$.
\item $\boxed{0}$.
\item $\boxed{-5\pi}$.
\item $\boxed{0}$.
\item $\boxed{0}$. The hemisphere together with its base disk is closed,
so its boundary is empty.
\item $\boxed{0}$. The entire cube surface is closed.
\end{enumerate}

\newpage
\begin{enumerate}[start=36]
\item Valid parametrizations and bounds are:
\begin{enumerate}[label=(\alph*),leftmargin=1.6em,itemsep=3pt]
\item $\vec r(u,v)=(u,v,10)$,
$-1\le u\le1$, $-2\le v\le2$.
The stated bounds describe a rectangle, rather than a square.
\item $\vec r(u,v)=(u,v,u^2+v^2)$,
$0\le v\le1$, $0\le u\le v$.
\item $\vec r(r,\theta)=(r\cos\theta,r\sin\theta,r^2)$,
$0\le r\le3$, $0\le\theta\le2\pi$.
\item $\vec r(r,\theta)=(r\cos\theta,r\sin\theta,r)$,
$2\le r\le4$, $0\le\theta\le2\pi$.
\item $\vec r(r,\theta)=(r\cos\theta,r\sin\theta,
4-r\cos\theta+r\sin\theta)$,
$0\le r\le3$, $0\le\theta\le2\pi$.
\item $\vec r(\theta,z)=(\cos\theta,\sin\theta,z)$,
$0\le\theta\le\pi/2$, $0\le z\le3$.
\item $\vec r(\theta,z)=(2\cos\theta,2\sin\theta,z)$,
$0\le\theta\le2\pi$,
$2\sin\theta-2\le z\le2\cos\theta+4$.
\item $\vec r(\phi,\theta)=(2\sin\phi\cos\theta,
2\sin\phi\sin\theta,2\cos\phi)$,
$0\le\phi\le\pi/4$, $0\le\theta\le2\pi$.
\end{enumerate}

\item $\boxed{\text{(a), (d), (e)}}$ under the course convention that
$\nabla$ denotes the gradient of a \emph{scalar} field; a Jacobian of
a vector field is not being requested in option (c).

\item $\boxed{\text{(a), (b), (d)}}$.
\item $\boxed{\text{(a), (b), (d)}}$.
\item In order: $\boxed{\mathrm{F},\mathrm{T},\mathrm{T},\mathrm{F},\mathrm{F},\mathrm{T}}$.

\item Choices are (a) direct parametrization, (b) FTLI, (c) Green,
(d) Stokes, (e) none. The coordinator's intended \emph{direct-method}
selections are listed below. The original question does not explicitly
exclude planar reductions; those also make Green valid in (3) and (5).
\begin{center}
\renewcommand{\arraystretch}{1.1}
\begin{tabular}{@{}clll@{}}
\toprule
Case&Direct selections&Allowing planar Green&Integral value\\
\midrule
(1)&(a), (c), (d)&same&$2\pi$\\
(2)&(a), (b), (c), (d)&same&$0$\\
(3)&(a), (d)&(a), (c), (d)&$0$\\
(4)&(a), (b)&same&$2$\\
(5)&(a), (d)&(a), (c), (d)&$-3/2$\\
(6)&(a), (b), (c), (d)&same&$0$\\
\bottomrule
\end{tabular}
\end{center}
In (2), (4), and (6), potentials are $x^2y$, $x^2y+xz$, and
$x^3y+y^2$, respectively. In (4) the curve is open. In (1) the planar
scalar curl is $2$. The three-dimensional curls in (3) and (5) are
$(1,1,1)$ and $(-1,-1,-1)$, respectively, so neither field is conservative.

For (3), substitute $z=y+2$, $\dd z=\dd y$ into the line integral:
\[
\vec F\cdot\dd\vec r=(y+2)\,\dd x+(x+y)\,\dd y,
\qquad Q_x-P_y=0.
\]
For (5), substitute $z=1-x-y$, $\dd z=-\dd x-\dd y$:
\[
\vec F\cdot\dd\vec r=(y-x)\,\dd x+(1-2x-y)\,\dd y,
\qquad Q_x-P_y=-3.
\]
Thus both curves are planar (in tilted planes), and Green can be applied
after these explicit reductions. The base triangle in (5) has area $1/2$.
\end{enumerate}
\end{document}
